\documentclass[11pt,a4paper]{article}
\usepackage[T1]{fontenc}
\usepackage{lmodern,microtype}
\usepackage[a4paper,margin=25mm]{geometry}
\usepackage{amsmath,amssymb,booktabs,array,tabularx}
\usepackage[dvipsnames]{xcolor}
\usepackage[colorlinks=true,linkcolor=MidnightBlue,urlcolor=MidnightBlue]{hyperref}
\usepackage{fancyhdr}
\pagestyle{fancy}
\fancyhf{}
\fancyhead[L]{\small Rank-five fusion categories}
\fancyhead[R]{\small Second revision round}
\fancyfoot[C]{\thepage}
\setlength{\headheight}{14pt}
\setlength{\emergencystretch}{2em}
\setlength{\parskip}{3pt}
\newcommand{\FPdim}{\operatorname{FPdim}}
\newcommand{\Tr}{\operatorname{Tr}}
\hypersetup{pdftitle={Rank-five classification: second revision round},pdfauthor={}}
\title{Rank-five classification\\[4pt]\large Further audit and revision letter\\Second revision round}
\author{}
\date{24 November 2026}
\begin{document}
\maketitle

\section{Outcome and baseline}

This letter records a fresh audit of the previously delivered
\path{Rank5_General_Sharp_Revised.zip}. It concerns changes relative to
that revised package, not a repetition of the first revision letter.
The earlier \path{Review_Report.tex/pdf} and
\path{Revision_Letter.tex/pdf} remain unchanged as historical records.

One new, definite arithmetic error was found in the quadratic descent
proof: a discrete lower bound omitted an allowed value. Its correction
still gives the required strict contradiction. The review also found a
hypothesis that should be stated more precisely: the rationality in the
twist-support lemma concerns the normalized $S$-matrix row. A weaker
reading in terms of the associated fusion character is false, and an
explicit counterexample is given below. Two further precision changes
specify an integral parameter and make source-version references
unambiguous.

The corrected argument retains all principal conclusions:
\[
\begin{gathered}
\text{every complex rank-five fusion category has a pseudo-unitary
Galois conjugate}\\[-2pt]
\text{and admits a spherical structure},\\
\mu\le4,\qquad \FPdim(\mathcal C)\le48+24\sqrt3,\\
\text{exactly seventeen based rings, each with a unitary realization}.
\end{gathered}
\]
No new faulty exclusion, omitted bounded ring, or incorrect model tensor
was found. The ordinary verification and both optional audits passed.
No existing search range, categorical filter, witness tensor, or model
tensor was changed.

The review covered all seventeen sections and the existing programs,
including the previous repairs. The descent and local arithmetic
sections were reassigned for independent scrutiny; the invertible-action,
Morita, small-generator, quintic, and endpoint arguments were also
reconstructed. This is a mathematical and computational audit, not a
formal proof-assistant certification or a guarantee that no further
error exists.

\section{Complete manuscript change ledger}

\subsection{R2-M01: the omitted value in Proposition 6.3}

\textbf{Location:} the final $M=9$ subcase of Proposition 6.3.
\par\noindent Source label: \path{prop:quadratic-small-conductor-120}.
The proof has established
\[
 s=1-\frac{\gamma}{4}-\frac J9>\frac14,
 \qquad \gamma\in\mathbb Z_{\ge1},\quad J\in\{2,4\}.
\]
It then claimed $s\ge11/36$. This overlooks $(\gamma,J)=(2,2)$,
which gives $s=5/18<11/36$. The complete list above $1/4$ is
\[
\begin{array}{c|c|c}
\gamma&J&s\\\hline
1&2&19/36\\
1&4&11/36\\
2&2&5/18
\end{array}
\]
The fourth choice with $\gamma\in\{1,2\}$ is below $1/4$;
every $\gamma\ge3$ is also excluded by that lower bound.

Both uses of $11/36$ in this inference have been replaced by $5/18$.
The revised proof identifies the minimizing pair and displays the exact
comparison:
\[
 2\left(\frac5{18}-\frac2{120}\right)^2+\frac{32}{81}
 =\frac{8609}{16200}
 =\frac12+\frac{509}{16200}>\frac12.
\]
Thus the upper bound $\max(E,E')\le120$ survives the repair. The
subsequent quadratic enumeration has the same finite range and the same
output. The new optional audit verifies the finite list and the strict
inequality with exact fractions.

After finding this omission, the neighboring small-unit calculations
were checked separately. At exponent one the three conductor cases
$5,8,12$ give no allowed pair. At exponent two the conductor-five case
gives only $\beta=\epsilon$; the conductor-eight case gives exactly
$(m,n)=(4,1),(6,2),(12,4),(20,7)$, with $J\in\{2,4\}$.
The deductions $M\le4$ and $9\le M\le10$, and every other numerical
comparison in the proposition, remain correct.

\subsection{R2-M02: rationality of the normalized modular row}

\textbf{Location:} Lemma 13.5, label \texttt{lem:odd-row}, and the
paragraph immediately following its proof.
The phrase ``a rational row'' now explicitly means
\[
                       S_{Ji}\in\mathbb Q\quad\text{for every }i,
\]
where $S$ is the normalized modular matrix. The associated formal
codegree is $q=S_{0J}^{-2}\in\mathbb Z_{>0}$.
The proof uses this hypothesis to make the diagonal Galois operator
fix the row with positive sign.

Rationality of the ratios $S_{Ji}/S_{J0}$ alone would not suffice.
For a concrete counterexample to that weaker reading, take
\[
 \mathcal D=\mathrm{SU}(2)_4^{\boxtimes3}
              \boxtimes\mathrm{Semion}^{\boxtimes2},
 \qquad J=(2,2,2,0,0).
\]
All dimensions are positive. The central charge is
$i^3(e^{\pi i/4})^2=1$, and $\theta_J=1$.
The character on each $\mathrm{SU}(2)_4$ factor is
$(1,0,-1,0,1)$ and that on each semion factor is $(1,1)$,
so the product character is integral. Nevertheless,
\[
 q=108,\qquad S_{J0}=\frac1{6\sqrt3},\qquad
 (ST^5S)_{J0}=-\frac1{6\sqrt3},\qquad (5,108)=1.
\]
Indeed, with $\omega=e^{2\pi i/3}$, the respective factor entries are
$(1-\omega^2)/3$ and $(1+i)/2$; cubing the first and squaring the
second gives the negative value above. The normalized row is irrational,
so the revised hypothesis excludes this example.

Every use in the rank-five manuscript meets the intended stronger
hypothesis. The added paragraph explains this directly: for the central
induction label of a rational character $\rho$ of the original fusion
ring,
\[
 S_{A_\rho,W}=\frac{\rho(F(W))}{f_\rho}\in\mathbb Q,
\]
because $\rho$ is integer-valued and its codegree is integral.
Applications to nonrational conjugate pairs instead use the separate
integral-block lemma. Consequently no application, twist restriction,
or certificate is weakened by this precision change.

\subsection{R2-M03: the parameter in Proposition 11.1 is integral}

The statement of Proposition 11.1 now specifies
$h\in\mathbb Z_{\ge0}$ in place of merely $h\ge0$.
Its proof already derives $k=6h$ with integral $h$, and later small-case
arguments use that discreteness. This states the established conclusion
explicitly; it changes neither the proof nor the allowed parameters.

\subsection{R2-M04: exact source-version locators}

The higher Gauss-sum formula used in Section 11 and in the quintic
classification in Section 12 was cited as ``Corollary 4.5'' of
Ng--Schopieray--Wang.
In the linked arXiv version, the required identity is Theorem 4.4,
equation (4.17), while 4.5 is an example. This does not establish that
the published article uses the same numbering.

Both in-text citations now explicitly refer to Theorem 4.4, equation
(4.17), in \href{https://arxiv.org/abs/1812.11234v2}{arXiv:1812.11234v2},
the version whose formula was checked. The bibliography now pins that
version, gives the journal issue, adds
\href{https://doi.org/10.1007/s00029-019-0499-2}{doi:10.1007/s00029-019-0499-2},
and states which numbering the manuscript uses. No mathematical input
has been replaced by a different result.

\section{R2-C01: a reproducible additional audit}

The sole code-file addition is
\path{final_census/verify_additional_audits.py}.
It supplements the existing verifier with the following targeted checks.

\paragraph{Discrete quadratic cutoff.}
It enumerates the complete admissible $(\gamma,J)$ list above and
verifies the corrected minimum and margin exactly.

\paragraph{Direct cyclotomic averages.}
For the conductor-twelve field $\mathbb Q(\sqrt3)$ used in the
higher-indicator certificates, the program computes Galois averages
directly in $\mathbb Q[X]/\Phi_M(X)$. It uses every relevant order
dividing $144$, every reduced numerator, and powers one through six.
The resulting \textbf{864 exact moments and 19 joint vectors} agree
with the packaged relative-trace formula. Keeping each vector intact
checks the correlations between different powers. The zero vector for
higher primary orders is included.

\paragraph{Restriction columns and Gram conditions.}
For all \textbf{14 induction or higher-indicator certificate rings},
the full Cartesian product of the allowed unit-induction columns agrees
with the optimized join. Direct exact principal-minor tests agree with
the native Gram pruning. For all \textbf{11 residual Gram matrices}
in the higher-indicator branches, native and direct Python enumeration
of further columns agree as well.

\paragraph{Coverage of the quartic generator ranges.}
For $d=A+Bt+Ct^2+Et^3$, the new check obtains the coefficient range for
$E$ by rational interval Lagrange interpolation, then ranges for
$C,B,A$ by divided differences and evaluation. It uses the same
defining field polynomials and certified root intervals, but none of
the original fixed coefficient cutoffs, fixed padding, or nested
neighbor-range construction. Its bounds enclose every integral vector
with distinguished value in $[2,M]$ and other values in
$[-3.17,3.17]$, an outward enlargement of $[-\sqrt{10},\sqrt{10}]$.
All these vectors occur in the packaged enumeration. Direct interval
evaluation also certifies every reported Perron upper endpoint.

\begin{center}
\begin{tabular}{rrrr}
\toprule
Conductor&Cap $M$&Interval candidates&Packaged candidates\\\midrule
15&125&934&935\\
15&155&1160&1162\\
20&153&858&860\\\bottomrule
\end{tabular}
\end{center}
These counts precede the nonprimitive filter. They do not replace the
manuscript's primitive-candidate counts $931,1158,856$.
The first-round upper-endpoint repair is thereby confirmed through a
separate range construction. Its error allowance
$15\sum_i|x_i|$ is conservative: at $|t|<2$, replacing the root by
its lower interval endpoint and rounding the scaled powers of degree
at most three changes each power by less than thirteen units.

\paragraph{Scope of independence.}
The new checks share the package's field data, root intervals, and
construction of individual restriction-column sets. They independently
check the trace calculation, column joins, Gram enumeration, and
coefficient-range coverage. They do not independently prove the
categorical hypotheses behind these computations. The new program
uses explicit exception-based checks; its imported assertion-dependent
generator rejects optimized Python, as verified in this round.

\section{Verification and further mathematical review}

All three final commands exited successfully:
\begin{verbatim}
python3 verify.py
python3 final_census/verify_independent_census.py
python3 final_census/verify_additional_audits.py
\end{verbatim}
The first command again passed all eighteen structural checks and
reproduced the complete census:
\begin{center}
\begin{tabular}{rrrrr}
\toprule
$\mu$&Generated&Within both bounds&Excluded&Realized\\\midrule
1&16&16&6&10\\
2&37&30&27&3\\
3&82&46&44&2\\
4&134&22&20&2\\\midrule
Total&269&114&97&17\\\bottomrule
\end{tabular}
\end{center}
The 97 exclusions remain partitioned as 76 $d$-number failures,
two subcategory-divisibility failures, three central-dimension norm
failures, two dimension-field failures, ten induction/second-moment
failures, and four higher-indicator certificates. All seventeen model
identifications pass.

The previous optional census audit again reconstructs the symmetry,
associativity, and necessary linear equations; its alternative replay
agrees on 220 self-dual rings and 49 non-self-dual rings, the latter
split into 47 one-pair and two two-pair cases. The separate Pieri
construction again agrees with the $\operatorname{Ad}(\mathrm{SU}(3)_4)$
model. As already documented, this replay shares the fraction-free
linear-completion routine and basis parametrization.

The mathematical review rechecked the following dependencies in
addition to the newly repaired paragraph:
\begin{itemize}
\item sphericalization and formal-codegree normalizations, common
dimension fields, simultaneous columns, and the choice of embedding;
\item both Hecke projections, the ordinary Fourier proposition, and
the quadratic, quartic, and quintic finite reductions;
\item the invertible-action cases, the $S_3$ untwisting, Morita
dimensions and induction pairings, conductor-eight saturation, and
the endpoint equivariantization;
\item the small-generator graph argument, quintic maximal-order
argument, local-order socle arguments, and generalized-indicator limit;
\item the quadratic profile eliminations, cubic and quartic dyadic
depth arguments, twist moments, and exact search boundaries.
\end{itemize}
The primary texts for the delicate induction, congruence, slope,
dimension-grading, near-group, and model-existence inputs were checked
where needed. No further confirmed error was found. Standard background
results are still cited; this review does not claim to have reproved
every external theorem. The recent dimension-block preprint is not
needed as an unverified theorem input because the manuscript contains
its own Fourier proof.

\section{Remaining package changes and traceability}

Besides R2-M01--R2-M04 and R2-C01, the complete delivery changes are:
\begin{enumerate}
\item The manuscript PDF was rebuilt from the corrected source. The
title, empty author fields, theorem numbering, seventeen-ring list,
and principal constants were retained.
\item This letter is new, in
\path{Revision_Letter_Round2.tex} and
\path{Revision_Letter_Round2.pdf}. Both earlier report pairs are
preserved unchanged.
\item The root \path{README.md} identifies the current letter,
distinguishes the previous reports and patches, and documents both
optional audit commands. The census README describes the added
program, its scope, and its file role.
\item The complete ordinary rerun is recorded in
\path{verification.txt}; its deterministic output agrees with the
previous transcript. The new
\path{additional_verification_round2.txt} records the actual output
of both optional reruns with their command labels. The previous
\path{independent_verification.txt} is preserved.
\item The new \path{revision_round2.patch} gives the exact changes
to existing text files relative to the preceding delivered package,
apart from the regenerated checksum manifest. Newly added files are
listed in this letter. The earlier \path{revision.patch} is preserved.
\item \path{SHA256SUMS.txt} was regenerated for the final package.
The archive and every distributed checksum were checked. Temporary
build files, compiled helpers, and Python caches are excluded.
\end{enumerate}

The two optional commands do not change the mathematical filters or
stored witnesses. Requirements remain Python with SymPy and NumPy,
and a C++17 compiler with OpenMP. The source uses standard LaTeX
packages. The final manuscript and letter were rendered and visually
checked before delivery.

For an exact baseline, the SHA-256 of the preceding delivered ZIP is
\begin{center}\scriptsize\ttfamily
4a6a53f4a8d4208391dcaf8681c8bc5fde8eed17fb512a0753f073e0102c3a73
\end{center}
and that of its \path{rank5.tex} is
\begin{center}\scriptsize\ttfamily
485c2982cec7d4ec46ce96faab6393162866d2ede5fffac1e04bae965988d9fd
\end{center}
Every change described here is measured against that exact revision.

\end{document}
